\documentclass[11pt]{article}
\usepackage[utf8]{inputenc}
\usepackage{amsmath}
\usepackage{amsfonts}
\usepackage{amssymb}
\usepackage{bm}
\usepackage{geometry}
\usepackage{xcolor}
\usepackage{booktabs}

\geometry{margin=0.8in}

\title{Sparse Poisson GLM with Off-Diagonal Memory Kernel:\\Model, Inference, and Training}
\author{Balewski, NERSC}
\date{\today}

\begin{document}

\maketitle

\begin{abstract}
Two-block inference scheme for a Poisson GLM with a lagged off-diagonal memory kernel on a non-stationary neuronal population. The model uses a diagonal lag-1 drive plus a shared history kernel on off-diagonal inputs. Alternating optimization yields tractable convex subproblems for the kernel (\(\lambda_2\) smoothness) and network parameters (\(\lambda_3\) off-diagonal sparsity).
\end{abstract}

\tableofcontents
\newpage

% ===============================================================
\section{Generative Model}
\label{sec:gen}
% ===============================================================

Population of \(N\) neurons, \(T\) time bins of width \(\Delta t\). Spike counts \(Y_t \in \mathbb{N}_0^{N}\), \(t=0,\dots,T-1\).

\paragraph{Parameters.}
Diagonal: \(A^{\mathrm{diag}} \in \mathbb{R}^{N}\) with \((A^{\mathrm{diag}})_i = (A_{\text{true}})_{ii}\).
Off-diagonal: \(A^{\mathrm{off}} \in \mathbb{R}^{N\times N}\), zeros on diagonal.
Kernel: \(\kappa \in \mathbb{R}^{M}\), 0-based indexing, anchor \(\kappa_0=1\), learnable entries \(\kappa_1,\dots,\kappa_{M-1}\).
Bias: \(B \in \mathbb{R}^{N}\).
Clipping threshold: \(\eta_{\mathrm{clip}}=5\).

\paragraph{Off-diagonal history.}
\[
S_t \;=\; \sum_{\ell=0}^{M-1} \kappa_\ell \, Y_{t-1-\ell} \;\in\; \mathbb{R}^{N},
\qquad t-1-\ell \ge 0.
\]

\paragraph{State equation.}
\[
\eta_{t,i} \;=\; (A^{\mathrm{diag}})_i \, Y_{t-1,i}
\;+\; \big(A^{\mathrm{off}} S_t\big)_i
\;+\; B_i,
\]
\[
\mu_{t,i} \;=\; \exp\!\big(\tilde{\eta}_{t,i}\big)\,\Delta t,\quad
\tilde{\eta}_{t,i} \;=\; \min\big(\eta_{t,i},\,\eta_{\mathrm{clip}}\big).
\]

\paragraph{Observation model.}
\[
Y_{t,i} \sim \mathrm{Poisson}(\mu_{t,i}), \quad i=1,\dots,N,\; t=0,\dots,T-1.
\]

% ===============================================================
\section{Objective and Identifiability}
\label{sec:objective}
% ===============================================================

\paragraph{Log-likelihood.}
Dropping the constant \(\log(Y_{t,i}!)\) terms:
\[
\mathcal{L} = \sum_{i=1}^N \sum_{t=0}^{T-1}
\left[ Y_{t,i}\,\tilde{\eta}_{t,i} - \exp(\tilde{\eta}_{t,i})\,\Delta t \right].
\]

\paragraph{Objective function.}
\begin{equation}
\mathcal{J}
\;=\;
\underbrace{\sum_{t=2}^{T}\sum_{i=1}^{N}
\left[\mu_{t,i}-Y_{t,i}\left(\tilde{\eta}_{t,i}+\log\Delta t\right)\right]}_{\text{Poisson NLL}}
\;+\;
\underbrace{\lambda_2 \sum_{\ell=2}^{M-1} \big(\kappa_\ell-\kappa_{\ell-1}\big)^2}_{\text{kernel smoothness (Block~1)}}
\;+\;
\underbrace{\lambda_3\,\overline{|A_{\mathrm{off}}|}}_{\text{off-diagonal sparsity (Block~2)}},
\label{eq:objective}
\end{equation}
where
\begin{equation}
\overline{|A_{\mathrm{off}}|}
\;=\;
\frac{1}{N(N-1)}
\sum_{\substack{i,j=1\\i\neq j}}^{N}|A_{ij}^{\mathrm{off}}|.
\label{eq:l1_offdiag}
\end{equation}

\paragraph{Identifiability.}
The product \(A^{\mathrm{off}}_{ij}\,\kappa_\ell\) has a scaling degeneracy, removed by the anchor \(\kappa_0=1\).
If clipping activates frequently, a censored Poisson likelihood should replace the standard one.

% ===============================================================
\section{Inference: Two-Block Alternating Scheme}
\label{sec:inference}
% ===============================================================

The bilinear structure \(A^{\mathrm{off}}_{ij}\,\kappa_\ell\) precludes joint convexity. We use block coordinate descent: each block is convex given the other fixed, guaranteeing a non-increasing objective sequence converging to a stationary point.

\subsection{Initialization}
Initialize \(A^{\mathrm{diag}}\) and \(B\) from mean firing rates and lag-1 covariances.
Initialize \(\kappa\) with a damped-oscillator profile or simply \(\kappa_0=1\) with small subsequent entries.

\subsection{Block~1 (kernel update)}
Fix \((A^{\mathrm{diag}}, A^{\mathrm{off}}, B)\). Update \(\kappa_1,\dots,\kappa_{M-1}\) by minimizing \(\mathcal{J}\) with the \(\lambda_2\) smoothness penalty. The problem is convex in \(\kappa\) since \(S_t\) is linear in \(\kappa\).

\subsection{Block~2 (network update)}
Fix \(\kappa\). Update \((A^{\mathrm{diag}}, A^{\mathrm{off}}, B)\) by minimizing the Poisson NLL plus the \(\lambda_3\) L1 penalty~\eqref{eq:l1_offdiag}. Dale's law sign constraints on columns of \(A^{\mathrm{off}}\) can be added as linear inequalities.

% ===============================================================
\section{Training Samples}
\label{sec:training}
% ===============================================================

Input: spike counts \(Y_t \in \mathbb{N}_0^{N}\), \(t=0,\dots,T-1\).
Let \(M_{\mathrm{cut}} \ge 2\) be the history cutoff.

\paragraph{Sample construction.}
Each sample at output time \(t_{\mathrm{out}} \in [M_{\mathrm{cut}},\, T-1]\):
\[
X^{(t_{\mathrm{out}})} = \big\{Y_{t_{\mathrm{out}}-1},\, Y_{t_{\mathrm{out}}-2},\, \dots,\, Y_{t_{\mathrm{out}}-M_{\mathrm{cut}}}\big\} \in \mathbb{N}_0^{N \times M_{\mathrm{cut}}},
\qquad
y^{(t_{\mathrm{out}})} = Y_{t_{\mathrm{out}}} \in \mathbb{N}_0^{N}.
\]

\paragraph{Sampling strategy.}
Draw \(t_{\mathrm{out}}\) densely from \([M_{\mathrm{cut}},\, T-1]\) with stride \(s\); optionally enforce a minimum separation to reduce autocorrelation. Total samples:
\[
N_{\mathrm{samples}} \approx \left\lfloor \frac{T - M_{\mathrm{cut}}}{s} \right\rfloor.
\]
For validation, sample windows within random non-overlapping blocks of length \(L \gg M_{\mathrm{cut}}\) to obtain a temporally separated hold-out set.

% ===============================================================
\section{Kernel Normalization}
\label{sec:lags}
% ===============================================================

We use 0-based lag indexing \(\ell \in \{0,\dots,M-1\}\) with anchor \(\kappa_0=1\). An alternative \(\|\kappa\|_1=1\) constraint was tested but significantly reduced non-stationarity in simulations, so the ground truth \(\kappa_{\mathrm{true}}\) does not obey it.

The \(\lambda_2\) smoothness penalty imposes no parametric form on \(\kappa\). If prior knowledge constrains the kernel shape, one can reparametrize \(\kappa\) with fewer parameters while keeping \(\kappa_0=1\).

% ===============================================================
\section{Implementation Notes}
\label{sec:practical}
% ===============================================================

\paragraph{Sparsity and sign constraints.}
The \(\lambda_3\) penalty~\eqref{eq:l1_offdiag} promotes sparsity; normalization by \(N(N-1)\) makes the penalty strength network-size invariant. Column sign constraints (Dale's law) can be imposed if neuron types are known.

\paragraph{Regularization tuning.}
Start \(\lambda_2\) small; increase if the kernel estimate oscillates. Tune \(\lambda_3\) to match expected connectivity density (e.g., 5--20\%).



\end{document}
